%%%%%%%%%%%%%%%%%%%%%%%%%%%%%%%%%%%%%%%%%%%%%%%%%%%%%%%%%%%%%%%%%%%%%%%%%%%%
\thanks

This work has been supported by NSERC-Canada grant \# ODGP0110050
and FCAR-Qu\'ebec grant \# 93ER1654 to the first author.
Fran\c cois Blouin, Marco Jacques, Fran\c cois Paradis, and
Jos\'ee Turgeon have participated in the design and implementation of 
{\tt LatMRG}, between 1988 and 1992.
We wish to thank L. Afflerbach and H. Grothe for helpful discussions (in 1989).

%%%%%%%%%%%%%%%%%%%%%%%%%%%%%%%%%%%%%%%%%%%%%%%%%%%%%%%%%%%%%%%%%%%%%%%%%%%%

\bibliography {random,simul}
\bibliographystyle {plain}

%%%%%%%%%%%%%%%%%%%%%%%%%%%%%%%%%%%%%%%%%%%%%%%%%%%%%%%%%%%%%%%%%%%%%%%%%%%%

\iffalse %%%%%%%%%%%%%%%
\begin {refer}

\bibitem {AFF86} Afflerbach, L.
The Sub-Lattice Structure of Linear Congruential Random Number Generators.
{\it Manuscripta Math.\/}, {\bf 55} (1986) 455--465.

\bibitem {AFF85} Afflerbach, L. and Grothe, H.
Calculation of Minkowski-Reduced Lattice Bases.
{\it Computing\/}, {\bf 35} (1985) 269--276.

\bibitem {AFF88a} Afflerbach, L. and Grothe, H.
The Lattice Structure of Pseudo-Random Vectors Generated by Matrix Generators.
{\it J. of Computational and Applied Math.\/}, {\bf 23} (1988) 127--131.

\bibitem {BRA87} Bratley, P., Fox, B. L. and Schrage, L. E.
{\it A Guide to Simulation\/}, second edition. Springer-Verlag, New York, 1987.

\bibitem {COU93} Couture, R. and L'Ecuyer, P.
``Computing the Shortest vector and a Minkowski Reduced Basis in a Lattice'',
In preparation.

\bibitem {COU93b} Couture, R. and L'Ecuyer, P.
``On the Lattice Structure of Certain Linear Congruential Sequences
Related to AWC/SWB Generators'',
{\em Mathematics of Computation\/}, To appear.

\bibitem {COU94} Couture, R. and L'Ecuyer, P.
``Orbits and Lattices for Linear Random Number Generators with 
Composite Moduli'', Submitted

\bibitem {DIE75} Dieter, U.  How to Calculate Shortest Vectors in a Lattice.
{\it Math. of Computation\/}, {\bf 29}, 131 (1975) 827--833.

%\bibitem {ECK85} Eckhardt, H., Koch, J., Mall, M. et Putfarken, P. (1985).
%{\it Logitech Modula-2 User's Manual\/},
%Logitech Inc., Redwood City, California.

\bibitem {EIC89} Eichenauer-Herrmann, J., Grothe, H. and Lehn, J.
On the Period Length of Pseudorandom Vector Sequences Generated by
Matrix Generators.
{\it Math. of Computation\/}, {\bf 52}, 185 (1989) 145--148.

\bibitem {FIN85} Fincke, U. and Pohst, M.
Improved Methods for Calculating Vectors of Short Length in a Lattice,
Including a Complexity Analysis.
{\em Mathematics of Computation}, {\bf 44}, 170 (1985) 463--471.

\bibitem {FIS86} Fishman, G. S. and Moore III, L. S.
An Exhaustive Analysis of
Multiplicative Congruential Random Number Generators with Modulus $2^{31}-1$.
{\it SIAM J. on Scientific and Statistical Computing\/} {\bf 7}, 1 (1986)
24--45.

\bibitem {FIS90} Fishman, G. S.
Multiplicative Congruential Random Number Generators with Modulus $2^\beta$:
An Exhaustive Analysis for $\beta=32$ and a Partial Analysis for $\beta=48$.
{\em Mathematics of Computation}, {\bf 54}, 189 (1990) 331--344.

\bibitem {GRO88} Grothe, H. Matrixgeneratoren zur Erzeugung gleichverteilter
Pseudozufallsvektoren (in german).
Dissertation (thesis), Tech. Hochschule Darmstadt, Germany, 1988.

\bibitem {GRU73} Grube, A.  Mehrfach rekursiv-erzeugte Pseudo-Zufallszahlen
(in german). {\it Zeitsch\-rift f\"ur angewandte Math. und
Mechanik\/} {\bf 53} (1973) T223--T225.

%\bibitem {KIN88} King, K. N.
%{\em TopSpeed Modula-2, User's Manual},
%Jensen \& Partners International, 1988.

\bibitem {KNU81} Knuth, D. E.
{\it The Art of Computer Programming~: Seminumerical
Algorithms\/}, vol. 2, second edition. Ad\-di\-son-Wesley, 1981.

\bibitem {LEC90} L'Ecuyer, P.,
``Random Numbers for Simulation'',
{\em Communications of the ACM\/}, {\bf 33}, 10 (1990), 85--97.

\bibitem {LEC94} L'Ecuyer, P.,
``Uniform Random Number Generation'',
{\em Annals of Operations Research\/} (1994), to appear.

\bibitem {LEC94b} L'Ecuyer, P.,
``Combined Multiple Recursive Generators'', Manuscript.

\bibitem {LEC93a} L'Ecuyer, P., Blouin, F., and Couture, R.,
``A Search for Good Multiple Recursive Random Number Generators'',
{\em ACM Transactions on Modeling and Computer Simulation\/},
{\bf 3}, 2 (1993), 87--98.

\bibitem {LEC93b} L'Ecuyer, P. and Couture, R.,
``An Implementation of the Lattice and Spectral Tests for Linear
Congruential Generators'', submitted for publication.

\bibitem {LEC88} L'Ecuyer, P., Perron, G. et Blouin, F.
``SENTIERS: Un logiciel Modula-2 pour l'arithm\'e\-ti\-que sur les grands
entiers'', Rapport no. DIUL-RT-8802, D\'ept. d'informatique, Univ. Laval, 1988.

\bibitem {LEC91} L'Ecuyer, P. and Tezuka, S.,
Structural Properties for Two Classes of Combined Random Number Generators.
{\em Mathematics of Computation\/} {\bf 57}, 196 (1991) 735--746.

\bibitem {LID86} Lidl, R. and Niederreiter, H.
{\it Introduction to Finite Fields and Their Applications\/}.
Cambridge University Press, Cambridge, 1986.

\bibitem {MAR68} Marsaglia, G. Random Numbers Fall Mainly in the Planes.
{\it Proceedings of the National Academy of Sciences of the United States of
America\/} {\bf 60} (1968) 25--28.

\bibitem {MOR88} Morain, F. Implementation of the Atkin-Gold\-was\-ser-Kilian
Primality Testing Algorithm. Rapport de recherche no. 911, INRIA, Rocquencourt,
France, 1988.

\bibitem {NIE86} Niederreiter, H.
A Pseudorandom Vector Generator Based on Finite Field Arithmetic.
{\it Math. Japonica\/}, {\bf 31}, 5 (1986) 759--774.

\bibitem {RIP83} Ripley, B. D. The Lattice Structure of Pseudo-random Number
Generators.  {\em Proc. Roy. Soc. London, Series A}, {\bf 389} (1983) 197--204.

\bibitem {TEZ93} Tezuka, S., L'Ecuyer, P., and Couture, R.
``On the Lattice Structure of the Add-with-carry and Subtract-with-borrow
Random Number Generators'',
{\em ACM Transactions of Modeling and Computer Simulation\/},
To appear.

\end {refer}
\fi  %%%%%%%%%%%%%%%%%%
